% Integratiemethodes. Zurbarán, Stilleven met vier vaten (ca. 1650): vier vaten, elk een omwentelingslichaam;
% hun inhoud en de oppervlakte van hun wand zijn integralen (hoofdstuk 3).
% Motief: een vaas als omwentelingslichaam rond een verticale as (x = 18), met profiel
% r(z) = 0,75 + 0,55 sin(1,1 z + 0,3) + 0,08 z voor z van 0 tot 5,4, op hoogte y = 4,3 + z.
% De cirkels staan in een lichte bovenaanzicht als ellipsen (verticale halve as 0,18 r). Eén smalle band tussen
% z = 1,8 en z = 2,2 is in fel: het stukje mantel 2 pi r(z) Delta s uit de formule voor de manteloppervlakte.
\def\omslagbestand{omslag-vaten.jpg}
\def\omslagbron{Francisco de Zurbarán, \textit{Stilleven met vier vaten}, ca. 1650 (publiek domein)}
\def\omslagrand{10.5}
\def\omslagmotief{%
  \randlijn
  \def\vr##1{(0.75+0.55*sin((1.1*##1+0.3) r)+0.08*##1)}%
  % de as
  \draw[white,line width=.6pt,dashed,opacity=.7] (18,3.9) -- (18,10.1);
  % parallelcirkels (alleen de voorkant)
  \foreach \z in {0.6,1.2,...,4.8}{%
    \pgfmathsetmacro{\r}{\vr{\z}}%
    \draw[white,line width=.5pt,opacity=.55] plot[domain=180:360,samples=40,smooth,variable=\t]
      ({18+\r*cos(\t)},{4.3+\z+0.18*\r*sin(\t)});}
  % de band in fel
  \pgfmathsetmacro{\ra}{\vr{1.8}}\pgfmathsetmacro{\rb}{\vr{2.2}}%
  \fill[fel] plot[domain=180:360,samples=40,smooth,variable=\t] ({18+\ra*cos(\t)},{6.1+0.18*\ra*sin(\t)})
    -- plot[domain=360:180,samples=40,smooth,variable=\t] ({18+\rb*cos(\t)},{6.5+0.18*\rb*sin(\t)}) -- cycle;
  % het profiel, links en rechts
  \draw[omslaglijn=1.2pt] plot[domain=0:5.4,samples=80,smooth,variable=\z] ({18-\vr{\z}},{4.3+\z});
  \draw[omslaglijn=1.2pt] plot[domain=0:5.4,samples=80,smooth,variable=\z] ({18+\vr{\z}},{4.3+\z});
  % bodem en opening
  \pgfmathsetmacro{\rn}{\vr{0}}\pgfmathsetmacro{\rt}{\vr{5.4}}%
  \draw[omslaglijn=1pt] (18,4.3) ellipse ({\rn} and {0.18*\rn});
  \draw[omslaglijn=1pt] (18,9.7) ellipse ({\rt} and {0.18*\rt});
}
